% TOP_PROBLEM: {"problemId": "problem.green-griffiths-lang-conjecture-on-entire-curves-in-varieties-of-general-type", "canonicalTitle": "Green-Griffiths-Lang Conjecture on entire curves in varieties of general type", "releaseRank": 96, "primaryDomain": "geometry_topology", "primaryDomainLabel": "Geometry & topology", "displayStatus": "reviewed_hold", "releaseStatus": "open"}
% !TeX program = lualatex
% ATTEMPT_STATUS: unresolved
\documentclass[11pt]{article}
\usepackage[margin=25mm]{geometry}
\usepackage{fontspec}
\setmainfont{Latin Modern Roman}
\usepackage{amsmath,amssymb,mathtools}
\usepackage{unicode-math}
\setmathfont{Latin Modern Math}
\usepackage{xurl,hyperref}
\hypersetup{colorlinks=true,urlcolor=blue}
\setcounter{secnumdepth}{0}
\setlength{\parindent}{0pt}
\setlength{\parskip}{5pt}
\setlength{\emergencystretch}{3em}
\begin{document}
\section*{\texorpdfstring{Green-Griffiths-Lang Conjecture on entire curves in varieties of general type}{Green-Griffiths-Lang Conjecture on entire curves in varieties of general type}}\label{96}
\subsection{Definitions and mathematical statement}
An entire curve is a nonconstant holomorphic map $f:{\mathbb{C}}\to X$. Algebraic degeneracy means $\overline{f({\mathbb{C}})}^{\rm Zar}\subsetneq X$. The catalogue's literal assertion quantifies a possibly different proper subvariety for each curve:
\[
 \forall f:{\mathbb{C}}\to X\text{ entire},\quad\exists Z_f\subsetneq X:\ f({\mathbb{C}})\subseteq Z_f.
\]
The stronger usual uniform Green--Griffiths--Lang formulation places \emph{all} such curves in one proper algebraic subset depending only on $X$; Demailly's Introduction states that formulation. The existence of jet-differential equations is an important method but not by itself equivalent to either degeneracy conclusion: one must control their common solution locus. These distinctions are recorded rather than silently repairing the source wording.
\par\smallskip\textit{Formulation and concept references:} \href{https://arxiv.org/abs/1011.3636}{[S1]}, \href{https://arxiv.org/html/1011.3636v2}{[E1]}.

\subsection{Short English statement}
Must every entire curve in a projective variety of general type lie inside a proper algebraic subset, with the stronger standard version asking for one subset covering them all?
\subsection{Research attempt}

\paragraph{Scope, process, and literature check.}
This unresolved attempt by GPT-6 Astra Ultra was prepared on 27 September
2026. Throughout the work, $X$ is smooth, projective and of general type
over $\mathbb C$, as intended by the catalogue title and English
statement. Both the literal, curve-dependent conclusion and the stronger
uniform conclusion are kept visible. The plan is to prove concrete
subcases and formulate an exact sufficient condition on jet equations;
the counterexample track tests stronger assertions that this condition
might tempt one to assume.

Demailly [S1, E1] proves the existence of negatively twisted jet
differentials and discusses the separate base-locus problem. That
distinction is the key external input here. The complement-of-plane-curves
setting of [S2] and the specialized constructions of [S3] are retained
as references, not used as theorems about every projective variety of
general type. For the elementary hyperbolic-curve input below, see
Demailly's account of uniformization in
\href{https://www-fourier.ujf-grenoble.fr/~demailly/manuscripts/viasm2012expanded.pdf}{[E2]}.

\paragraph{A complete subcase and a product extension.}
Let $C$ be a smooth projective curve of genus at least two. Its universal
cover is the unit disk $\mathbb D$. Every holomorphic map
$f:\mathbb C\to C$ lifts to a holomorphic map
$\widetilde f:\mathbb C\to\mathbb D$, since $\mathbb C$ is simply
connected. The lift is bounded and hence constant by Liouville's theorem.
Thus there are no entire curves in $C$; the stronger uniform assertion
holds with empty exceptional set.

For $Y=C_1\times\cdots\times C_r$ with all genera at least two,
compose a map $f:\mathbb C\to Y$ with each projection. Each component
is constant, so $f$ is constant. These examples meet the general-type
hypothesis: the external tensor product of the ample canonical bundles
of the curves is $K_Y$, and is ample. This argument also applies to
any closed subvariety of the product as a statement about its entire
maps, whether or not that subvariety is itself of general type.

There is a useful extension with a different conclusion. If a variety
$X$ admits a surjective morphism $\pi:X\to C$ to a curve of genus at
least two and $\dim X>1$, then every entire curve lies in one fibre:
$\pi\circ f$ is constant by the preceding argument. The fibre is a
proper algebraic subset because $\pi$ is surjective and nonconstant.
This proves the literal curve-dependent assertion for such an $X$,
including all general-type examples with this morphism. It does not
produce one proper subset containing the union of the relevant fibres.
Additional information about which fibres contain entire curves is
needed for that conclusion.

\paragraph{General type does not mean absence of entire curves.}
Take $Y=C_1\times C_2$ as above and let $X$ be its blow-up at a point.
If $\pi:X\to Y$ denotes the blow-down, the exceptional divisor is
$E\simeq\mathbb P^1$. The canonical divisor satisfies
$K_X=\pi^*K_Y+E$. The pullback of the big divisor $K_Y$ is big, and
adding the effective divisor $E$ preserves bigness. Thus $X$ is of
general type. The standard map $z\mapsto[z:1]$ followed by
$E\hookrightarrow X$ is a nonconstant entire curve.

Nevertheless the stronger conjecture holds here with $Z=E$. Indeed
$\pi\circ f$ is a map to the hyperbolic product and is constant. Over
any point other than the blown-up point, the fibre of $\pi$ is a single
point, so a nonconstant $f$ must lie in $E$. This example verifies both
the necessity of allowing an exceptional subset and the difference
between algebraic degeneracy and hyperbolicity. It is a test of the
overstrong claim, not a counterexample to the catalogue statement.

\paragraph{The jet-equation strategy with all quantifiers exposed.}
Fix a positive integer $k$. A $k$-jet records the derivatives through
order $k$ of a germ $(\mathbb C,0)\to X$; call it regular if its first
derivative is nonzero. In local coordinates a jet differential is a
polynomial in these derivatives, with weight $j$ assigned to the
$j$th derivative and with the prescribed coordinate transformation law.
For sections with suitable negative ample twists, the fundamental
vanishing theorem says that evaluation on any entire curve vanishes.
We use that theorem as an external analytic input, not as a consequence
of the mere existence of a polynomial.

Here is the exact additional condition that would complete this route.
Suppose finitely many such global sections $P_1,\ldots,P_s$, all viewed
at order $k$, and a nonempty Zariski open subset $U\subset X$ satisfy
\begin{equation}\label{r96:separation}
 \text{for every }x\in U\text{ and every regular jet }\xi\text{ over }x,
 \quad P_j(\xi)\ne0\text{ for some }j.
\end{equation}
The twist and weight may differ between the sections; nonvanishing at
a jet is meaningful in each one-dimensional value space.

\paragraph{Lemma 96.1: jet separation gives uniform degeneracy.}
Under the vanishing theorem and \eqref{r96:separation}, every entire
curve has image in $X\setminus U$.

\emph{Proof.} Suppose a nonconstant $f$ meets $U$. The set
$f^{-1}(U)$ is nonempty and open in $\mathbb C$. The zeros of $f'$ are
discrete unless $f$ is constant: in a coordinate neighbourhood, the
derivatives of its component functions are holomorphic and cannot all
vanish on an open set without making the map locally constant.
Consequently there exists $t\in f^{-1}(U)$ where $f'(t)\ne0$.
Its regular jet is annihilated by every $P_j$, by the vanishing theorem,
contradicting \eqref{r96:separation}.\hfill$\square$

The proof handles critical points without asserting that all jets of a
nonconstant curve are regular. It also supplies a single exceptional
set determined before choosing $f$, so it proves more than the literal
quantifier statement whenever its assumptions hold.

\paragraph{A fibre calculation that breaks a shortcut.}
At first order, take one homogeneous differential of positive degree
$m$ on a variety of dimension $n\ge2$. At a point $x$ its value is a
homogeneous polynomial $P_x(v)$ on $T_xX\simeq\mathbb C^n$, after
trivializing its line-bundle value. It always has a nonzero zero. If it
is identically zero there is nothing to show. Otherwise restrict it to
a two-dimensional subspace on which it is not zero. A nonzero binary
homogeneous polynomial of positive degree has a root in $\mathbb P^1$
over $\mathbb C$, hence a nonzero zero in that subspace.

Thus even a section nonzero at a generic point as a tensor need not
exclude all nonzero tangent directions at any point. Its zero set in
the projectivized tangent bundle can project onto the whole base.
For example $P(v_1,v_2)=v_1^m$ leaves the direction $[0:1]$ everywhere.
Every tangent vector is the derivative of some local germ, so this is
a genuine defect of a proposed first-order jet-separation argument.
It does not assert that those germs extend to entire curves.

Two equations can behave differently: $v_1^m$ and $v_2^m$ have no
common nonzero zero in dimension two. The hard geometric issue is to
obtain enough global negatively twisted sections with this joint
behaviour, not simply to increase a vector-space dimension count.
The toy polynomials are not claimed to globalize on an arbitrary $X$.

At higher order, the zero jet also annihilates every positive-weight
polynomial at every base point. Therefore a base-locus argument using
the entire affine jet bundle, including zero jets, cannot have proper
projection at all. One must remove the unsuitable jets or work in the
appropriate projectivized construction and still control its regular
part. Removing that elementary obstruction alone does not prove
\eqref{r96:separation}.

\paragraph{Why the weaker quantifier is not automatically uniform.}
The inference from ``each curve is in some proper closed subset'' to
``one proper closed subset contains all curves'' is logically invalid.
For a simple algebraic set-theoretic model, the lines in $\mathbb P^2$
are each proper and their union is the whole plane. This is not offered
as a general-type example satisfying the weak conjecture; it only tests
the quantifier interchange. Noetherianity controls descending chains
of closed sets or ascending chains of ideals, not arbitrary unions of
proper closed subsets.

Similarly, finite generation of an ideal of equations at a fixed jet
order says that its zero locus has a finite description. It does not
say that the zero locus has proper image in $X$. Nor does it furnish
a uniform order for a process that continually introduces new jet
variables. Those are separate tasks in a proposed inductive proof.

\paragraph{Verification, first gap, and next targets.}
The disk-lifting and product arguments prove their stated subcases.
The blow-up has big canonical divisor and an explicit entire curve,
but its image stays in the proper exceptional divisor. The jet lemma
uses universal vanishing and regular-jet separation as separate
assumptions; neither is inferred from a raw count of sections. The
two-variable polynomial checks are exact algebra, and finite sampling
of jets or entire maps would not verify the universal condition.

The first missing statement for the chosen route is
\eqref{r96:separation}, or a weaker geometric condition forcing every
entire jet trajectory in the common zero locus to project into a
proper subset. The next tasks are to compute the regular common zero
locus for a specific higher-dimensional family, test restriction maps
of jet sections on its components dominating $X$, and study whether
the surviving directions admit global entire integral curves. A
counterexample to the literal conjecture must instead be a
Zariski-dense entire curve on a general-type variety; the exceptional
curve above plainly fails that requirement.

\paragraph{Final status.}
\textbf{UNRESOLVED.} Complete subcase arguments and the conditional
jet-separation lemma are established. The passage from existence of
jet equations to control of all their entire solutions remains open
in this attempt. None of the elementary tests is claimed as new.

\subsection{Sources}
\begin{itemize}
\item[S1]J.-P. Demailly, Holomorphic Morse inequalities and the Green-Griffiths-Lang conjecture, arXiv:1011.3636 [math.AG], 2010.
\url{https://arxiv.org/abs/1011.3636}.
\textit{Formulation source.}
\item[S2]Computing Jet Differentials and the Green-Griffiths-Lang Conjecture for Complements of Smooth Plane Curves.
\url{https://arxiv.org/abs/2608.20479}.
\textit{Further reference; not a proof endorsement.}
\item[S3]Non-reductive geometric invariant theory and hyperbolicity.
\url{https://link.springer.com/article/10.1007/s00222-023-01219-z}.
\textit{Further reference; not a proof endorsement.}
\item[E1]Jean-Pierre Demailly, Holomorphic Morse inequalities and the Green-Griffiths-Lang conjecture, arXiv:1011.3636v2 (2010), Introduction, Conjecture (0.2) and discussion of jet differentials.
\url{https://arxiv.org/html/1011.3636v2}.
\textit{Added primary source: Quantifier distinction and limitations of the jet-differential reformulation. Consulted 24 September 2026.}
\item[E2]J.-P. Demailly, Hyperbolic algebraic varieties and holomorphic differential equations, expanded lecture notes, author-hosted version.
\url{https://www-fourier.ujf-grenoble.fr/~demailly/manuscripts/viasm2012expanded.pdf}.
\textit{Uniformization and hyperbolic-curve background; consulted 27 September 2026.}
\end{itemize}
\end{document}
